\section*{Bab \ref{ch:b2:anthropogenic-climate-change} --- Perubahan Iklim Antropogenik}

\begin{solution}{exo:b2:anthropogenic-climate-change:1}
(a) $3.75 - 0.60 + 0.22 + 0.45 + 0.11 = \qty{3.93}{W.m^{-2}}$, dengan penyesuaian-penyesuaian yang menambahkan
\qty{0.18}{W.m^{-2}}. (b) $3.93/3.71 = 1.059$: \qty{5.9}{\percent} lebih besar. (c) Penyesuaian suhu
troposfer, $-0.60$: troposfer yang lebih hangat memancarkan lebih banyak inframerah ke angkasa, yang membuang
sebagian energi tambahan itu sebelum permukaan menanggapinya.
\end{solution}

\begin{solution}{exo:b2:anthropogenic-climate-change:2}
(a) $4\pi(\num{6.371e6})^2 = \qty{5.10e14}{m^2}$. (b) Selama 47 tahun, $\num{434.9e21}/(47 \times
\num{3.156e7} \times \num{5.10e14}) = \qty{0.575}{W.m^{-2}}$; selama 12 tahun, $\num{152.4e21}/(12
\times \num{3.156e7} \times \num{5.10e14}) = \qty{0.789}{W.m^{-2}}$, yaitu 0.57 dan 0.79 dalam
penilaian itu. (c) $0.91 \times 434.9 = \qty{396}{ZJ}$.
\end{solution}

\begin{solution}{exo:b2:anthropogenic-climate-change:3}
(a) $-3.22 + 1.30 + 0.35 + 0.42 - 0.01 = \num{-1.16}$, yaitu nilai bersihnya. (b) $3.93/1.16 =
\qty{3.39}{\celsius}$; $3.93/3.22 = \qty{1.22}{\celsius}$. (c) $3.39/1.22 = 2.78$: umpan balik itu
hampir melipattigakan pemanasan Planck.
\end{solution}

\begin{solution}{exo:b2:anthropogenic-climate-change:4}
(a) Satu petagram karbon setara dengan $44.009/12.011 = 3.664$~Gt karbon dioksida: $0.45 \times 3.664 =
\qty{1.65}{\celsius}$ per \qty{1000}{PgC}, dan rentangnya menjadi 0.989 sampai 2.31. (b) $2390 \times
12.011/44.009 = \qty{652}{PgC}$. (c) $0.45 \times 2.39 = \qty{1.08}{\celsius}$.
\end{solution}

\begin{solution}{exo:b2:anthropogenic-climate-change:5}
(a) $g = 159/398 = 0.400$; $T = (398/\num{5.670e-8})^{1/4} = \qty{289}{K}$. (b) $-0.6 \times 4
\times \num{5.670e-8} \times 288^3 = \qty{-3.25}{W.m^{-2}.K^{-1}}$; pada \qty{298}{K}, $-3.60$.
(c) Kurang peka: tanggapan itu tumbuh sebagai $T^3$, sehingga planet yang lebih hangat lebih mudah melepaskan setiap derajat
tambahan, bila umpan balik lainnya sama.
\end{solution}

\begin{solution}{exo:b2:anthropogenic-climate-change:6}
(a) $1.85 - 0.50 = \qty{1.35}{W.m^{-2}.\celsius^{-1}}$. (b) $\sqrt{0.32^2 + 0.20^2} = 0.38$.
(c) Jumlahnya jauh lebih baik diketahui daripada masing-masing suku: galat-galatnya berkorelasi negatif dengan kuat. Di tempat
troposfer atas menghangat lebih banyak, troposfer itu juga menampung lebih banyak uap, sehingga \omterm{def:b2:anthropogenic-climate-change:feedbacks}{umpan balik laju lapse} yang lebih
negatif datang bersama \omterm{def:b2:anthropogenic-climate-change:feedbacks}{umpan balik uap air} yang lebih positif.
\end{solution}

\begin{solution}{exo:b2:anthropogenic-climate-change:7}
(a) $\lambda = (0.59 - 2.20)/1.03 = \qty{-1.56}{W.m^{-2}.\celsius^{-1}}$; $3.93/1.56 =
\qty{2.52}{\celsius}$. (b) $3.93/(1.56 - 0.5) = 3.93/1.06 = \qty{3.71}{\celsius}$. (c) $3.93/1.46 =
\qty{2.69}{\celsius}$. Penilaian itu, dengan distribusi lengkap, mendapatkan 3.5 dan \qty{2.7}{\celsius}.
\end{solution}

\begin{solution}{exo:b2:anthropogenic-climate-change:8}
(a) $0.617/3682 = \qty{1.68e-4}{\celsius^{-1}}$. (b) $0.113 \times 0.396 = \qty{0.0447}{m}$, atau
\qty{44.7}{mm}, dekat dengan \qty{47.5}{mm} dalam anggaran itu. (c) $\num{2e-4} \times 0.2 \times 700 =
\qty{0.028}{m}$, \qty{28}{mm}.
\end{solution}

\begin{solution}{exo:b2:anthropogenic-climate-change:9}
(a) $\sum xy = 0.311$, $\sum x^2 = 0.300$: $\beta = 1.04$. (b) Selang \qty{95}{\percent}-nya adalah
$1.04 \pm 1.96 \times 0.08$, 0.88 sampai 1.20: selang itu tidak memuat 0, sehingga tanggapan itu terdeteksi, dan memuat
1, sehingga perubahan yang teramati konsisten dengan yang disimulasikan. (c) Selang 0.34 sampai 0.66 masih
tidak memuat 0 (terdeteksi), tetapi tidak memuat 1: model itu akan melebih-lebihkan tanggapan dengan faktor sekitar
dua.
\end{solution}

\begin{solution}{exo:b2:anthropogenic-climate-change:10}
(a) $\sum g^2 = 1.20$, $\sum a^2 = 0.18$, $\sum ga = -0.46$, $\sum gy = 0.74$, $\sum ay = -0.28$.
Persamaan normalnya adalah $1.20\,\beta_g - 0.46\,\beta_a = 0.74$ dan $-0.46\,\beta_g + 0.18\,\beta_a
= -0.28$; determinannya $1.20 \times 0.18 - 0.46^2 = 0.0044$, dan $\beta_g = (0.74 \times
0.18 - 0.46 \times 0.28)/0.0044 = 1.00$, $\beta_a = (-1.20 \times 0.28 + 0.46 \times 0.74)/0.0044 =
1.00$. (b) Gas rumah kaca $+\qty{0.80}{\celsius}$, aerosol $-\qty{0.30}{\celsius}$. (c) Determinan itu
kecil karena $g$ dan $a$ hampir sebanding (korelasinya $-0.99$): sedikit
derau dalam $y$ dapat diserap dengan menaikkan kedua faktor skala bersama-sama, sehingga masing-masing kurang
terkekang walaupun efek gabungannya diketahui dengan baik.
\end{solution}

\begin{solution}{exo:b2:anthropogenic-climate-change:11}
(a) $(0.93 - 0.1)/0.45 \times 1000 = \qty{1844}{Gt}$; dikurangi 210 dan $26 \times 0.93 = 24$:
\qty{1610}{Gt}. (b) $0.93 - x = 0.45 \times (1350 + 210 + 24)/1000 = 0.713$, sehingga $x =
\qty{0.22}{\celsius}$: kira-kira ujung atas rentang hasil penilaian 0.1 sampai \qty{0.2}{\celsius}. (c) Anggaran
itu bergantung pada pemanasan non-CO$_2$ $x$ hanya melalui $-x/\text{TCRE}$: setiap
\qty{0.1}{\celsius} mengurangi $0.1/0.45 \times 1000 = \qty{222}{Gt}$, berapa pun batasnya.
\end{solution}

\begin{solution}{exo:b2:anthropogenic-climate-change:12}
(a) $3.93/1.81 = \qty{2.17}{\celsius}$, $3.93/0.51 = \qty{7.71}{\celsius}$, $3.93/1.16 =
\qty{3.39}{\celsius}$. (b) Sisi bawah sepanjang $3.39 - 2.17 = \qty{1.22}{\celsius}$, sisi atas
$7.71 - 3.39 = \qty{4.32}{\celsius}$: fungsi $1/|\lambda|$ cembung, sehingga langkah-langkah $\lambda$ yang sama
ke arah nol makin lama makin menaikkan ECS. (c) $\lambda$ yang lemah negatif memberikan ECS yang sangat
besar, dan perubahan kecil $\lambda$ di dekat nol sulit disingkirkan; ECS yang rendah memerlukan $\lambda$ yang sangat
negatif, yang dibatasi oleh \omterm{def:b2:anthropogenic-climate-change:feedbacks}{respons Planck} dan dibantah oleh pemanasan yang teramati.
\end{solution}

\begin{solution}{pb:b2:anthropogenic-climate-change:1}
\textbf{1.} $2.16 + 0.54 + 0.21 + 0.41 = \qty{3.32}{W.m^{-2}}$.
\textbf{2.} $-0.22 - 0.84 = \qty{-1.06}{W.m^{-2}}$.
\textbf{3.} $3.32 + 0.47 + 0.05 - 1.06 - 0.20 + 0.08 + 0.06 = \qty{2.72}{W.m^{-2}}$, yaitu total
hasil penilaian.
\textbf{4.} $3.93\log_2(410/278) = \qty{2.20}{W.m^{-2}}$, dalam batas \qty{2}{\percent} dari 2.16.
\textbf{5.} $3.93\log_2(427.35/278) = \qty{2.44}{W.m^{-2}}$.
\textbf{6.} Suku aerosol--awan, $-1.45$ sampai $-0.25$, selebar \qty{1.20}{W.m^{-2}}, dibandingkan dengan 1.52 untuk
total: suku itu membawa sebagian besar ketidakpastian total karena besar dan kurang
terkekang.
\textbf{7.} $0.79 \times \num{5.10e14} \times 12 \times \num{3.156e7} = \qty{1.53e23}{J}$,
\qty{153}{ZJ}.
\textbf{8.} $0.91 \times 153 = \qty{139}{ZJ}$.
\textbf{9.} $\lambda = (0.59 - 2.20)/1.03 = \qty{-1.56}{W.m^{-2}.\celsius^{-1}}$.
\textbf{10.} $3.93/1.56 = \qty{2.52}{\celsius}$.
\textbf{11.} $3.93 \times 1.03/2.20 = \qty{1.84}{\celsius}$.
\textbf{12.} $3.93/1.16 = \qty{3.39}{\celsius}$, lebih tinggi \qty{0.87}{\celsius}.
\textbf{13.} Pemanasan sejauh ini paling kuat di tempat umpan balik paling negatif; seiring berkembangnya
pola itu, parameter umpan balik menjadi kurang negatif, sehingga sensitivitas jangka panjang
melebihi sensitivitas efektif.
\textbf{14.} $0.45 \times 44.009/12.011 = \qty{1.65}{\celsius}$ per \qty{1000}{PgC}.
\textbf{15.} 0.989 sampai \qty{2.31}{\celsius} per \qty{1000}{PgC}.
\textbf{16.} $2390 \times 12.011/44.009 = \qty{652}{PgC}$.
\textbf{17.} $0.45 \times 2.39 = \qty{1.08}{\celsius}$.
\textbf{18.} Dalam batas pembulatan, karbon dioksida saja menjelaskan seluruh pemanasan manusia:
pemanasan oleh gas rumah kaca lain dan pendinginan oleh aerosol hampir saling meniadakan.
\textbf{19.} $1.5 - 1.07 = \qty{0.43}{\celsius}$.
\textbf{20.} $(0.43 - 0.1)/0.45 \times 1000 = \qty{733}{Gt}$ karbon dioksida.
\textbf{21.} $733 - 210 = \qty{523}{Gt}$.
\textbf{22.} $26 \times 0.43 = \qty{11.2}{Gt}$; $523 - 11.2 = \qty{512}{Gt}$.
\textbf{23.} Anggaran itu turun sebesar $0.1/0.45 \times 1000 = \qty{222}{Gt}$, menjadi \qty{290}{Gt}: yaitu
\qty{220}{Gt} yang menurut penilaian itu merupakan besar pergeseran anggaran akibat pilihan tentang pendorong-pendorong lain.
\textbf{24.} Bila dihitung dari TCRE, \omterm{def:b2:anthropogenic-climate-change:budget}{sisa anggaran karbon} untuk \qty{1.5}{\celsius} sejak
awal 2020 adalah \textbf{\qty[text-series-to-math=true]{512}{Gt}} karbon dioksida, dalam batas
\qty{3}{\percent} dari \qty{500}{Gt} hasil penilaian pada probabilitas \qty{50}{\percent}.
\end{solution}
